\Chapter{8}

\Verse{1}
{
    \zA{话说宝玉和凤姐回家，见过众人，宝玉便回明贾母要约秦钟上家塾之事，自己 也有个伴读的朋友，正好发愤； 又着实称赞秦钟人品行事，最是可人怜爱的。 凤姐
        又在一旁帮着说：“改日秦钟还来拜见老祖宗呢。” 说的贾母喜欢起来。 凤姐又趁 势请贾母一同过去看戏。 贾母虽年高，却极有兴头。
        后日，尤氏来请，遂带了王夫 人、黛玉、宝玉等过去看戏。 至晌午，贾母便回来歇息。 王夫人本好清净，见贾母 回来，也就回来了。
        然后凤姐坐了首席，尽欢至晚而罢。}
}
{
    \eA{Pao-yü and lady Feng, we will now explain, paid, on their return home, their
        respects to all the inmates, and Pao-yü availed himself of the first
        occasion to tell dowager lady Chia of his wish that Ch'in Chung should come
        over to the family school.}
}
{
}

\Verse{2}
{
    \zA{却说宝玉送贾母回来，待贾母歇了中觉，还要回去看戏，又恐搅的秦氏等人不 便。 因想起宝钗近日在家养病，未去看视，意欲去望他。 若从上房后角门过去，恐
        怕遇见别事缠绕，又怕遇见他父亲，更为不妥，宁可绕个远儿。 当下众嬷嬷丫鬟伺 候他换衣服，见不曾换，仍出二门去了，众嬷嬷丫鬟只得跟随出来。 还只当他去那
        边府中看戏，谁知到了穿堂儿，便向东北边绕过厅后而去。}
}
{
    \eA{"The presence for himself of a friend as schoolmate would," he argued, "be
        fitly excellent to stir him to zeal," and he went on to speak in terms of
        high praise of Ch'in Chung, his character and his manners, which most of all
        made people esteem him. Lady Feng besides stood by him and backed his
        request. "In a day or two," she added, "Ch'in Chung will be coming to pay
        his obeisance to your venerable ladyship."}
}
{
}

\Verse{3}
{
    \zA{偏顶头遇见了门下清客 相公詹光、单聘仁二人走来，一见了宝玉，便都赶上来笑着，一个抱着腰，一个拉
        着手，道：“我的菩萨哥儿!我说做了好梦呢，好容易遇见你了！” 说着，又唠叨 了半日才走开。 老嬷嬷叫住，因问：“你们二位是往老爷那里去的不是？” 二人点
        头道：“是。” 又笑着说：“老爷在梦坡斋小书房里歇中觉呢，不妨事的。” 一面 说，一面走了，说的宝玉也笑了。 于是转弯向北奔梨香院来。}
}
{
    \eA{This bit of news greatly rejoiced the heart of dowager lady Chia, and lady
        Feng likewise did not let the opportunity slip, without inviting the old
        lady to attend the theatrical performance to come off the day after the
        morrow.}
}
{
}

\Verse{4}
{
    \zA{可巧管库房的总领吴 新登和仓上的头目名叫戴良的，同着几个管事的头目，共七个人从帐房里出来，一 见宝玉，赶忙都一齐垂手站立。
        独有一个买办名唤钱华，因他多日未见宝玉，忙上 来打千儿请宝玉的安，宝玉含笑伸手叫他起来。 众人都笑说：“前儿在一处看见二
        爷写的斗方儿，越发好了，多早晚赏我们几张贴贴。” 宝玉笑道：“在那里看见了？” 众人道：“好几处都有，都称赞的了不得，还和我们寻呢！”
        宝玉笑道：“不值什 么，你们说给我的小么儿们就是了。”}
}
{
    \eA{Dowager lady Chia was, it is true, well on in years, but was, nevertheless,
        very fond of enjoyment, so that when the day arrived and Mrs. Yu came over
        to invite her round, she forthwith took madame Wang, Lin Tai-yü, Pao-yü and
        others along and went to the play. It was about noon, when dowager lady Chia
        returned to her apartments for her siesta;}
}
{
}

\Verse{5}
{
    \zA{一面说，一面前走，众人待他过去，方都各 自散了。 闲言少述。 且说宝玉来至梨香院中，先进薛姨妈屋里来，见薛姨妈打点针黹与 丫鬟们呢。
        宝玉忙请了安，薛姨妈一把拉住，抱入怀中笑说：“这么冷天，我的儿， 难为你想着来!快上炕来坐着罢。” 命人沏滚滚的茶来。 宝玉因问：“哥哥没在家 么？”
        薛姨妈叹道：“他是没笼头的马，天天逛不了，那里肯在家一日呢？” 宝玉 道：“姐姐可大安了？” 薛姨妈道：“可是呢，你前儿又想着打发人来瞧他。}
}
{
    \eA{and madame Wang, who was habitually partial to a quiet life, also took her
        departure after she had seen the old lady retire. Lady Feng subsequently
        took the seat of honour; and the party enjoyed themselves immensely till the
        evening, when they broke up. But to return to Pao-yü.}
}
{
}

\Verse{6}
{
    \zA{他在 里间不是，你去瞧。 他那里比这里暖和，你那里坐着，我收拾收拾就进来和你说话 儿。” 宝玉听了，忙下炕来到了里间门前，只见吊着半旧的红绸软帘。
        宝玉掀帘一步 进去，先就看见宝钗坐在炕上作针线，头上挽着黑漆油光的？ 儿，蜜合色的棉袄，
        玫瑰紫二色金银线的坎肩儿，葱黄绫子棉裙：一色儿半新不旧的，看去不见奢华， 惟觉雅淡。 罕言寡语，人谓装愚； 安分随时，自云守拙。
        宝玉一面看，一面问：“姐 姐可大愈了？”}
}
{
    \eA{Having accompanied his grandmother Chia back home, and waited till her
        ladyship was in her midday sleep, he had in fact an inclination to return to
        the performance, but he was afraid lest he should be a burden to Mrs. Ch'in
        and the rest and lest they should not feel at ease.}
}
{
}

\Verse{7}
{
    \zA{宝钗抬头看见宝玉进来，连忙起身含笑答道：“已经大好了，多谢 惦记着。” 说着，让他在炕沿上坐下，即令莺儿：“倒茶来。” 一面又问老太太姨
        娘安，又问别的姐妹们好。 一面看宝玉头上戴着累丝嵌宝紫金冠，额上勒着二龙捧 珠抹额，身上穿着秋香色立蟒白狐腋箭袖，系着五色蝴蝶鸾绦，项上挂着长命锁、
        记名符，另外有那一块落草时衔下来的宝玉。 宝钗因笑说道：“成日家说你的这块 玉，究竟未曾细细的赏鉴过，我今儿倒要瞧瞧。” 说着便挪近前来。}
}
{
    \eA{Remembering therefore that Pao Ch'ai had been at home unwell for the last
        few days, and that he had not been to see her, he was anxious to go and look
        her up, but he dreaded that if he went by the side gate, at the back of the
        drawing-room, he would be prevented by something or other, and fearing, what
        would be making matters worse, lest he should come across his father, he
        consequently thought it better to go on his way by a detour.}
}
{
}

\Verse{8}
{
    \zA{宝玉亦凑过去， 便从项上摘下来，递在宝钗手内。 宝钗托在掌上，只见大如雀卵，灿若明霞，莹润 如酥，五色花纹缠护。
        看官们须知道，这就是大荒山中青埂峰下的那块顽石幻相。 后人有诗嘲云： 女娲炼石已荒唐，又向荒唐演大荒。 失去本来真面目，幻来新就臭皮囊。
        好知运败金无彩，堪叹时乖玉不光。 白骨如山忘姓氏，无非公子与红妆。 通灵宝玉正面　　　　　　　通灵宝玉反面
        那顽石亦曾记下他这幻相并癞僧所镌篆文，今亦按图画于后面。}
}
{
    \eA{The nurses and waiting-maids thereupon came to help him to change his
        clothes; but they saw him not change, but go out again by the second door.
        These nurses and maids could not help following him out; but they were still
        under the impression that he was going over to the other mansion to see the
        theatricals.}
}
{
}

\Verse{9}
{
    \zA{但其真体最小，方 从胎中小儿口中衔下，今若按式画出，恐字迹过于微细，使观者大废眼光，亦非畅 事，所以略展放些，以便灯下醉中可阅。
        今注明此故，方不至以胎中之儿口有多大、 怎得衔此狼？ 蠢大之物为诮。 宝钗看毕，又从新翻过正面来细看，口里念道：“莫失莫忘，仙寿恒昌。” 念
        了两遍，乃回头向莺儿笑道：“你不去倒茶，也在这里发呆作什么？” 莺儿也嘻嘻 的笑道：“我听这两句话，倒像和姑娘项圈上的两句话是一对儿。”}
}
{
    \eA{Contrary to their speculations, upon reaching the entrance hall, he
        forthwith went to the east, then turned to the north, and walking round by
        the rear of the hall, he happened to come face to face with two of the
        family companions, Mr. Ch'an Kuang, and Mr. Tan T'ing-jen.}
}
{
}

\Verse{10}
{
    \zA{宝玉听了，忙 笑道：“原来姐姐那项圈上也有字？ 我也赏鉴赏鉴。” 宝钗道：“你别听他的话， 没有什么字。” 宝玉央及道：“好姐姐，你怎么瞧我的呢！”
        宝钗被他缠不过，因 说道：“也是个人给了两句吉利话儿，錾上了，所以天天带着。 不然沉甸甸的，有 什么趣儿？”
        一面说，一面解了排扣，从里面大红袄儿上将那珠宝晶莹、黄金灿烂 的璎珞摘出来。 宝玉忙托着锁看时，果然一面有四个字，两面八个字，共成两句吉 谶。
        亦曾按式画下形相。}
}
{
    \eA{As soon as they caught sight of Pao-yü, they both readily drew up to him,
        and as they smiled, the one put his arm round his waist, while the other
        grasped him by the hand. "Oh divine brother!" they both exclaimed, "this we
        call dreaming a pleasant dream, for it's no easy thing to come across you!"
        While continuing their remarks they paid their salutations, and inquired
        after his health;}
}
{
}

\Verse{11}
{
    \zA{金锁正面　　　　　　金锁反面 宝玉看了，也念了两遍，又念自己的两遍，因笑问：“姐姐，这八个字倒和我的是 一对儿。”
        莺儿笑道：“是个癞头和尚送的，他说必须錾在金器上――”宝钗不等 他说完，便嗔着：“不去倒茶！” 一面又问宝玉从那里来。
        宝玉此时与宝钗挨肩坐着，只闻一阵阵的香气，不知何味，遂问：“姐姐熏的 是什么香？ 我竟没闻过这味儿。” 宝钗道：“我最怕熏香。 好好儿的衣裳，为什么
        熏他？” 宝玉道：“那么着这是什么香呢？”}
}
{
    \eA{and it was only after they had chatted for ever so long, that they went on
        their way. The nurse called out to them and stopped them, "Have you two
        gentlemen," she said, "come out from seeing master?" They both nodded
        assent. "Your master," they explained, "is in the Meng P'o Chai small
        library having his siesta; so that you can go through there with no fear."
        As they uttered these words, they walked away.}
}
{
}

\Verse{12}
{
    \zA{宝钗想了想，说：“是了，是我早起 吃了冷香丸的香气。” 宝玉笑道：“什么‘冷香丸’，这么好闻？ 好姐姐，给我一 丸尝尝呢。” 宝钗笑道：“又混闹了。
        一个药也是混吃的？” 一语未了，忽听外面人说：“林姑娘来了。” 话犹未完，黛玉已摇摇摆摆的进 来，一见宝玉，便笑道：“哎哟!我来的不巧了。”
        宝玉等忙起身让坐。 宝钗笑道： “这是怎么说？” 黛玉道：“早知他来，我就不来了。” 宝钗道：“这是什么意思？”}
}
{
    \eA{This remark also evoked a smile from Pao-yü, but without further delay he
        turned a corner, went towards the north, and came into the Pear Fragrance
        Court, where, as luck would have it, he met the head manager of the
        Household Treasury, Wu Hsin-teng, who, in company with the head of the
        granary, Tai Liang, and several other head stewards, seven persons in all,
        was issuing out of the Account Room.}
}
{
}

\Verse{13}
{
    \zA{黛玉道：“什么意思呢：来呢一齐来，不来一个也不来； 今儿他来，明儿我来，间 错开了来，岂不天天有人来呢？ 也不至太冷落，也不至太热闹。 姐姐有什么不解的
        呢？” 宝玉因见他外面罩着大红羽缎对襟褂子，便问：“下雪了么？” 地下老婆们 说：“下了这半日了。” 宝玉道：“取了我的斗篷来。”
        黛玉便笑道：“是不是？ 我来了他就该走了！” 宝玉道：“我何曾说要去，不过拿来预备着。”}
}
{
    \eA{On seeing Pao-yü approaching, they, in a body, stood still, and hung down
        their arms against their sides. One of them alone, a certain butler, called
        Ch'ien Hua, promptly came forward, as he had not seen Pao-yü for many a day,
        and bending on one knee, paid his respects to Pao-yü. Pao-yü at once gave a
        smile and pulled him up.}
}
{
}

\Verse{14}
{
    \zA{宝玉的奶母 李嬷嬷便说道：“天又下雪，也要看时候儿，就在这里和姐姐妹妹一处玩玩儿罢。 姨太太那里摆茶呢。 我叫丫头去取了斗篷来，说给小么儿们散了罢？”
        宝玉点头。 李嬷嬷出去，命小厮们：“都散了罢。” 这里薛姨妈已摆了几样细巧茶食，留他们喝茶吃果子。 宝玉因夸前日在东府里 珍大嫂子的好鹅掌。
        薛姨妈连忙把自己糟的取了来给他尝。 宝玉笑道：“这个就酒 才好！” 薛姨妈便命人灌了上等酒来。 李嬷嬷上来道：“姨太太，酒倒罢了。”}
}
{
    \eA{"The day before yesterday," smiled all the bystanders, "we were somewhere
        together and saw some characters written by you, master Secundus, in the
        composite style. The writing is certainly better than it was before! When
        will you give us a few sheets to stick on the wall?" "Where did you see
        them?" inquired Pao-yü, with a grin.}
}
{
}

\Verse{15}
{
    \zA{宝 玉笑央道：“好妈妈，我只喝一钟。” 李妈道：“不中用，当着老太太、太太，那 怕你喝一坛呢。
        不是那日我眼错不见，不知那个没调教的只图讨你的喜欢，给了你 一口酒喝，葬送的我挨了两天骂!姨太太不知道他的性子呢，喝了酒更弄性。 有一
        天老太太高兴，又尽着他喝； 什么日子又不许他喝。 何苦我白赔在里头呢？” 薛姨 妈笑道：“老货!只管放心喝你的去罢。 我也不许他喝多了。
        就是老太太问，有我 呢！”}
}
{
    \eA{"They are to be found in more than one place," they replied, "and every one
        praises them very much, and what's more, asks us for a few." "They are not
        worth having," observed Pao-yü smilingly; "but if you do want any, tell my
        young servants and it will be all right." As he said these words, he moved
        onwards. The whole party waited till he had gone by, before they separated,
        each one to go his own way.}
}
{
}

\Verse{16}
{
    \zA{一面命小丫头：“来，让你奶奶去也吃一杯搪搪寒气。” 那李妈听如此说， 只得且和众人吃酒去。 这里宝玉又说：“不必烫暖了，我只爱喝冷的。” 薛姨妈道：
        “这可使不得：吃了冷酒，写字手打颤儿。” 宝钗笑道：“宝兄弟，亏你每日家杂 学旁收的，难道就不知道酒性最热，要热吃下去，发散的就快； 要冷吃下去，便凝
        结在内。 拿五脏去暖他，岂不受害？ 从此还不改了呢。 快别吃那冷的了。” 宝玉听 这话有理，便放下冷的，令人烫来方饮。}
}
{
    \eA{But we need not dilate upon matters of no moment, but return to Pao-yü. On
        coming to the Pear Fragrance Court, he entered, first, into "aunt" Hsüeh's
        room, where he found her getting some needlework ready to give to the
        waiting-maids to work at. Pao-yü forthwith paid his respects to her, and
        "aunt" Hsüeh, taking him by the hand, drew him towards her and clasped him
        in her embrace.}
}
{
}

\Verse{17}
{
    \zA{黛玉磕着瓜子儿，只管抿着嘴儿笑。 可巧黛玉的丫鬟雪雁走来给黛玉送小手炉 儿，黛玉因含笑问他说：“谁叫你送来的？ 难为他费心。 那里就冷死我了呢！” 雪
        雁道：“紫鹃姐姐怕姑娘冷，叫我送来的。” 黛玉接了，抱在怀中，笑道：“也亏 了你倒听他的话!我平日和你说的，全当耳旁风，怎么他说了你就依，比圣旨还快
        呢。” 宝玉听这话，知是黛玉借此奚落，也无回复之词，只嘻嘻的笑了一阵罢了。 宝钗素知黛玉是如此惯了的，也不理他。}
}
{
    \eA{"With this cold weather," she smilingly urged, "it's too kind of you, my
        dear child, to think of coming to see me; come along on the stove-couch at
        once!--Bring some tea," she continued, addressing the servants, "and make it
        as hot as it can be!" "Isn't Hsüeh P'an at home?" Pao-yü having inquired:
        "He's like a horse without a halter," Mrs. Hsüeh remarked with a sigh;}
}
{
}

\Verse{18}
{
    \zA{薛姨妈因笑道：“你素日身子单弱，禁不 得冷，他们惦记着你倒不好？” 黛玉笑道：“姨妈不知道：幸亏是姨妈这里，倘或 在别人家，那不叫人家恼吗？
        难道人家连个手炉也没有，巴巴儿的打家里送了来？ 不 说丫头们太小心，还只当我素日是这么轻狂惯了的呢。” 薛姨妈道：“你是个多心 的，有这些想头。
        我就没有这些心。” 说话时，宝玉已是三杯过去了，李嬷嬷又上来拦阻。 宝玉正在个心甜意洽之时， 又兼姐妹们说说笑笑，那里肯不吃？}
}
{
    \eA{"he's daily running here and there and everywhere, and nothing can induce
        him to stay at home one single day." "Is sister (Pao Ch'ai) all right
        again?" asked Pao-yü. "Yes," replied Mrs. Hsüeh, "she's well again. It was
        very kind of you two days ago to again think of her, and send round to
        inquire after her. She's now in there, and you can go and see her. It's
        warmer there than it's here;}
}
{
}

\Verse{19}
{
    \zA{只得屈意央告：“好妈妈，我再吃两杯就不吃 了。” 李嬷嬷道：“你可仔细今儿老爷在家，提防着问你的书！” 宝玉听了此话，
        便心中大不悦，慢慢的放下酒，垂了头。 黛玉忙说道：“别扫大家的兴。 舅舅若叫， 只说姨妈这里留住你。 这妈妈，他又该拿我们来醒脾了！”
        一面悄悄的推宝玉，叫 他赌赌气，一面咕哝说：“别理那老货，咱们只管乐咱们的。” 那李妈也素知黛玉 的为人，说道：“林姐儿，你别助着他了。
        你要劝他只怕他还听些。”}
}
{
    \eA{go and sit with her inside, and, as soon as I've put everything away, I'll
        come and join you and have a chat." Pao-yü, upon hearing this, jumped down
        with alacrity from the stove-couch, and walked up to the door of the inner
        room, where he saw hanging a portière somewhat the worse for use, made of
        red silk.}
}
{
}

\Verse{20}
{
    \zA{黛玉冷笑道： “我为什么助着他？ ――我也不犯着劝他。 你这妈妈太小心了!往常老太太又给他酒 吃，如今在姨妈这里多吃了一口，想来也不妨事。
        必定姨妈这里是外人，不当在这 里吃，也未可知。” 李嬷嬷听了，又是急，又是笑，说道：“真真这林姐儿，说出 一句话来，比刀子还利害。”
        宝钗也忍不住笑着把黛玉腮上一拧，说道：“真真的 这个颦丫头一张嘴，叫人恨又不是，喜欢又不是。”}
}
{
    \eA{Pao-yü raised the portière and making one step towards the interior, he
        found Pao Ch'ai seated on the couch, busy over some needlework. On the top
        of her head was gathered, and made into a knot, her chevelure, black as
        lacquer, and glossy like pomade. She wore a honey-coloured wadded robe, a
        rose-brown short-sleeved jacket, lined with the fur of the squirrel of two
        colours: the "gold and silver;"}
}
{
}

\Verse{21}
{
    \zA{薛姨妈一面笑着，又说：“别 怕，别怕，我的儿!来到这里没好的给你吃，别把这点子东西吓的存在心里，倒叫 我不安。 只管放心吃，有我呢!索性吃了晚饭去。
        要醉了，就跟着我睡罢。” 因命： “再烫些酒来。 姨妈陪你吃两杯，可就吃饭罢。” 宝玉听了，方又鼓起兴来。 李嬷
        嬷因吩咐小丫头：“你们在这里小心着，我家去换了衣裳就来。” 悄悄的回薛姨妈 道：“姨太太别由他尽着吃了。” 说着便家去了。}
}
{
    \eA{and a jupe of leek-yellow silk. Her whole costume was neither too new,
        neither too old, and displayed no sign of extravagance. Her lips, though not
        rouged, were naturally red; her eyebrows, though not pencilled, were yet
        blue black; her face resembled a silver basin, and her eyes, juicy plums.
        She was sparing in her words, chary in her talk, so much so that people said
        that she posed as a simpleton.}
}
{
}

\Verse{22}
{
    \zA{这里虽还有两三个老婆子，都是不关痛痒的，见李妈走了，也都悄悄的自寻方 便去了。 只剩了两个小丫头，乐得讨宝玉的喜欢。 幸而薛姨妈千哄万哄，只容他吃
        了几杯，就忙收过了。 作了酸笋鸡皮汤，宝玉痛喝了几碗，又吃了半碗多碧粳粥； 一时薛林二人也吃完了饭，又酽酽的喝了几碗茶。 薛姨妈才放了心。
        雪雁等几个人， 也吃了饭进来伺候。 黛玉因问宝玉道：“你走不走？” 宝玉乜斜倦眼道：“你要走 我和你同走。”}
}
{
    \eA{She was quiet in the acquittal of her duties and scrupulous as to the proper
        season for everything. "I practise simplicity," she would say of herself.
        "How are you? are you quite well again, sister?" inquired Pao-yü, as he
        gazed at her; whereupon Pao Ch'ai raised her head, and perceiving Pao-yü
        walk in, she got up at once and replied with a smile, "I'm all right again;}
}
{
}

\Verse{23}
{
    \zA{黛玉听说，遂起身道：“咱们来了这一日，也该回去了。” 说着， 二人便告辞。 小丫头忙捧过斗笠来，宝玉把头略低一低，叫他戴上。 那丫头便将这
        大红猩毡斗笠一抖，才往宝玉头上一合，宝玉便说：“罢了罢了!好蠢东西，你也 轻些儿。 难道没见别人戴过？ 等我自己戴罢。” 黛玉站在炕沿上道：“过来，我给
        你戴罢。” 宝玉忙近前来。 黛玉用手轻轻笼住束发冠儿，将笠沿掖在抹额之上，把 那一颗核桃大的绛绒簪缨扶起，颤巍巍露于笠外。}
}
{
    \eA{many thanks for your kindness in thinking of me." While uttering this, she
        pressed him to take a seat on the stove-couch, and as he sat down on the
        very edge of the couch, she told Ying Erh to bring tea and asked likewise
        after dowager lady Chia and lady Feng. "And are all the rest of the young
        ladies quite well?" she inquired.}
}
{
}

\Verse{24}
{
    \zA{整理已毕，端详了一会，说道： “好了，披上斗篷罢。” 宝玉听了，方接了斗篷披上。 薛姨妈忙道：“跟你们的妈 妈都还没来呢，且略等等儿。”
        宝玉道：“我们倒等着他们!有丫头们跟着就是了。” 薛姨妈不放心，吩咐两个女人送了他兄妹们去。 他二人道了扰，一径回至贾母房中。
        贾母尚未用晚饭，知是薛姨妈处来，更加 喜欢。 因见宝玉吃了酒，遂叫他自回房中歇着，不许再出来了。 又令人好生招呼着。
        忽想起跟宝玉的人来，遂问众人：“李奶子怎么不见？”}
}
{
    \eA{Saying this she scrutinised Pao-yü, who she saw had a head-dress of
        purplish-gold twisted threads, studded with precious stones. His forehead
        was bound with a gold circlet, representing two dragons, clasping a pearl.
        On his person he wore a light yellow, archery-sleeved jacket, ornamented
        with rampant dragons, and lined with fur from the ribs of the silver fox;}
}
{
}

\Verse{25}
{
    \zA{众人不敢直说他家去了， 只说：“才进来了，想是有事，又出去了。” 宝玉踉跄着回头道：“他比老太太还 受用呢，问他作什么!没有他只怕我还多活两日儿。”
        一面说，一面来至自己卧室。 只见笔墨在案。 晴雯先接出来，笑道：“好啊，叫我研了墨，早起高兴，只写了三 个字，扔下笔就走了，哄我等了这一天。
        快来给我写完了这些墨才算呢！” 宝玉方 想起早起的事来，因笑道：“我写的那三个字在那里呢？” 晴雯笑道：“这个人可 醉了。}
}
{
    \eA{and was clasped with a dark sash, embroidered with different-coloured
        butterflies and birds. Round his neck was hung an amulet, consisting of a
        clasp of longevity, a talisman of recorded name, and, in addition to these,
        the precious jade which he had had in his mouth at the time of his birth.}
}
{
}

\Verse{26}
{
    \zA{你头里过那府里去，嘱咐我贴在门斗儿上的。 我恐怕别人贴坏了，亲自爬高 上梯，贴了半天，这会子还冻的手僵着呢！” 宝玉笑道：“我忘了。 你手冷，我替
        你渥着。” 便伸手拉着晴雯的手，同看门斗上新写的三个字。 一时黛玉来了，宝玉笑道：“好妹妹，你别撒谎，你看这三个字那一个好？”
        黛玉仰头看见是“绛芸轩”三字，笑道：“个个都好，怎么写的这样好了!明儿也 替我写个匾。” 宝玉笑道：“你又哄我了。” 说着又问：“袭人姐姐呢？”}
}
{
    \eA{"I've daily heard every one speak of this jade," said Pao Ch'ai with a
        smile, "but haven't, after all, had an opportunity of looking at it closely,
        but anyhow to-day I must see it." As she spoke, she drew near. Pao-yü
        himself approached, and taking it from his neck, he placed it in Pao Ch'ai's
        hand. Pao Ch'ai held it in her palm.}
}
{
}

\Verse{27}
{
    \zA{晴雯向 里间炕上努嘴儿。 宝玉看时，见袭人和衣睡着。 宝玉笑道：“好啊!这么早就睡了。” 又问晴雯道：“今儿我那边吃早饭，有一碟子豆腐皮儿的包子。
        我想着你爱吃，和 珍大奶奶要了，只说我晚上吃，叫人送来的。 你可见了没有？” 晴雯道：“快别提 了。 一送来我就知道是我的。 偏才吃了饭，就搁在那里。
        后来李奶奶来了看见，说： ‘宝玉未必吃了，拿去给我孙子吃罢。’ 就叫人送了家去了。” 正说着，茜雪捧上 茶来。 宝玉还让：“林妹妹喝茶。”}
}
{
    \eA{It appeared to her very much like the egg of a bird, resplendent as it was
        like a bright russet cloud; shiny and smooth like variegated curd and
        covered with a net for the sake of protection. Readers, you should know that
        this was the very block of useless stone which had been on the Ta Huang
        Hills, and which had dropped into the Ch'ing Keng cave, in a state of
        metamorphosis.}
}
{
}

\Verse{28}
{
    \zA{众人笑道：“林姑娘早走了，还让呢。” 宝玉 吃了半盏，忽又想起早晨的茶来，问茜雪道：“早起沏了碗枫露茶，我说过那茶是
        三四次后才出色，这会子怎么又斟上这个茶来？” 茜雪道：“我原留着来着，那会 子李奶奶来了，喝了去了。” 宝玉听了，将手中茶杯顺手往地下一摔，豁琅一声打
        了个粉碎，泼了茜雪一裙子。 又跳起来问着茜雪道：“他是你那一门子的‘奶奶’， 你们这么孝敬他？
        不过是我小时候儿吃过他几日奶罢了，如今惯的比祖宗还大!撵出 去大家干净！”}
}
{
    \eA{A later writer expresses his feelings in a satirical way as follows:  Nü
        Wo's fusion of stones was e'er a myth inane,  But from this myth hath sprung
        fiction still more insane! Lost is the subtle life, divine, and real!--gone!
        Assumed, mean subterfuge! foul bags of skin and bone! Fortune, when once
        adverse, how true! gold glows no more! In evil days, alas! the jade's
        splendour is o'er!}
}
{
}

\Verse{29}
{
    \zA{说着立刻便要去回贾母。 原来袭人未睡，不过是故意儿装睡，引着宝玉来怄他玩耍。 先听见说字问包子， 也还可以不必起来；
        后来摔了茶钟动了气，遂连忙起来解劝。 早有贾母那边的人来 问：“是怎么了？” 袭人忙道：“我才倒茶，叫雪滑倒了，失手砸了钟子了。” 一
        面又劝宝玉道：“你诚心要撵他也好，我们都愿意出去，不如就势儿连我们一齐撵 了，你也不愁没有好的来伏侍你。” 宝玉听了，方才不言语了。}
}
{
    \eA{Bones, white and bleached, in nameless hill-like mounds are flung,  Bones
        once of youths renowned and maidens fair and young. The rejected stone has
        in fact already given a record of the circumstances of its transformation,
        and the inscription in seal characters, engraved upon it by the bald-headed
        bonze, and below will now be also appended a faithful representation of it;}
}
{
}

\Verse{30}
{
    \zA{袭人等便搀至炕上， 脱了衣裳，不知宝玉口内还说些什么，只觉口齿缠绵，眉眼愈加饧涩，忙伏侍他睡 下。 袭人摘下那“通灵宝玉”来，用绢子包好，？
        在褥子底下，恐怕次日带时冰了 他的脖子。 那宝玉到枕就睡着了。 彼时李嬷嬷等已进来了，听见醉了，也就不敢上 前，只悄悄的打听睡着了，方放心散去。
        次日醒来，就有人回：“那边小蓉大爷带了秦钟来拜。” 宝玉忙接出去，领了 拜见贾母。}
}
{
    \eA{but its real size is so very diminutive, as to allow of its being held by a
        child in his mouth while yet unborn, that were it to have been drawn in its
        exact proportions, the characters would, it is feared, have been so
        insignificant in size, that the beholder would have had to waste much of his
        eyesight, and it would besides have been no pleasant thing.}
}
{
}

\Verse{31}
{
    \zA{贾母见秦钟形容标致，举止温柔，堪陪宝玉读书，心中十分喜欢，便留 茶留饭，又叫人带去见王夫人等。 众人因爱秦氏，见了秦钟是这样人品，也都欢喜，
        临去时都有表礼。 贾母又给了一个荷包和一个金魁星，取“文星和合”之意。 又嘱 咐他道：“你家住的远，或一时冷热不便，只管住在我们这里。 只和你宝二叔在一
        处，别跟着那不长进的东西们学。” 秦钟一一的答应，回家禀知他父亲。}
}
{
    \eA{While therefore its shape has been adhered to, its size has unavoidably been
        slightly enlarged, to admit of the reader being able, conveniently, to
        peruse the inscription, even by very lamplight, and though he may be under
        the influence of wine.}
}
{
}

\Verse{32}
{
    \zA{他父亲秦邦业现任营缮司郎中，年近七旬，夫人早亡，因年至五旬时尚无儿女， 便向养生堂抱了一个儿子和一个女儿。 谁知儿子又死了，只剩下个女儿，小名叫做
        可儿，又起个官名叫做兼美。 长大时，生得形容袅娜，性格风流，因素与贾家有些 瓜葛，故结了亲。 秦邦业却于五十三岁上得了秦钟，今年十二岁了； 因去岁业师回
        南，在家温习旧课，正要与贾亲家商议附往他家塾中去。}
}
{
    \eA{These explanations have been given to obviate any such sneering remarks as:
        "What could be, pray, the size of the mouth of a child in his mother's womb,
        and how could it grasp such a large and clumsy thing?" On the face of the
        jade was written:  Precious Gem of Spiritual Perception. If thou wilt lose
        me not and never forget me,  Eternal life and constant luck will be with
        thee!}
}
{
}

\Verse{33}
{
    \zA{可巧遇见宝玉这个机会， 又知贾家塾中司塾的乃现今之老儒贾代儒，秦钟此去，可望学业进益，从此成名， 因十分喜悦。
        只是宦囊羞涩，那边都是一双富贵眼睛：少了拿不出来。 因是儿子的 终身大事所关，说不得东并西凑，恭恭敬敬封了二十四两贽见礼，带了秦钟到代儒
        家来拜见，然后听宝玉拣的好日子一同入塾。 塾中从此闹起事来。 未知如何，下回分解。 (本章完)}
}
{
    \eA{On the reverse was written:  1 To exorcise evil spirits and the accessory
        visitations; 2 To cure predestined sickness; 3 To prognosticate weal and
        woe. Pao Ch'ai having looked at the amulet, twisted it again to the face,
        and scrutinising it closely, read aloud:  If thou wilt lose me not and never
        forget me,  Eternal life and constant luck will be with thee!}
}
{
}

\Verse{34}
{
    \zA{}
}
{
    \eA{She perused these lines twice, and, turning round, she asked Ying Erh
        laughingly: "Why don't you go and pour the tea? what are you standing here
        like an idiot!" "These two lines which I've heard," smiled Ying Erh, "would
        appear to pair with the two lines on your necklet, miss!" "What!" eagerly
        observed Pao-yü with a grin, when he caught these words, "are there really
        eight characters too on your necklet, cousin? do let me too see it." "Don't
        listen to what she says," remarked Pao Ch'ai, "there are no characters on
        it." "My dear cousin," pleaded Pao-yü entreatingly, "how is it you've seen
        mine?" Pao Ch'ai was brought quite at bay by this remark of his, and she
        consequently added, "There are also two propitious phrases engraved on this
        charm, and that's why I wear it every day.}
}
{
}

\Verse{35}
{
    \zA{}
}
{
    \eA{Otherwise, what pleasure would there be in carrying a clumsy thing." As she
        spoke, she unfastened the button, and produced from inside her crimson robe,
        a crystal-like locket, set with pearls and gems, and with a brilliant golden
        fringe. Pao-yü promptly received it from her, and upon minute examination,
        found that there were in fact four characters on each side; the eight
        characters on both sides forming two sentences of good omen. The similitude
        of the locket is likewise then given below. On the face of the locket is
        written:  "Part not from me and cast me not away;" And on the reverse:  "And
        youth, perennial freshness will display!"}
}
{
}

\Verse{36}
{
    \zA{}
}
{
    \eA{Pao-yü examined the charm, and having also read the inscription twice over
        aloud, and then twice again to himself, he said as he smiled, "Dear cousin,
        these eight characters of yours form together with mine an antithetical
        verse." "They were presented to her," ventured Ying Erh, "by a mangy-pated
        bonze, who explained that they should be engraved on a golden trinket...."
        Pao Ch'ai left her no time to finish what she wished to say, but speedily
        called her to task for not going to bring the tea, and then inquired of
        Pao-yü "Where he had come from?" Pao-yü had, by this time, drawn quite close
        to Pao Ch'ai, and perceived whiff after whiff of some perfume or other, of
        what kind he could not tell. "What perfume have you used, my cousin," he
        forthwith asked, "to fumigate your dresses with?}
}
{
}

\Verse{37}
{
    \zA{}
}
{
    \eA{I really don't remember smelling any perfumery of the kind before." "I'm
        very averse," replied Pao Ch'ai blandly, "to the odour of fumigation; good
        clothes become impregnated with the smell of smoke." "In that case,"
        observed Pao-yü, "what scent is it?" "Yes, I remember," Pao Ch'ai answered,
        after some reflection; "it's the scent of the 'cold fragrance' pills which I
        took this morning." "What are these cold fragrance pills," remarked Pao-yü
        smiling, "that they have such a fine smell? Give me, cousin, a pill to try."
        "Here you are with your nonsense again," Pao Ch'ai rejoined laughingly; "is
        a pill a thing to be taken recklessly?" She had scarcely finished speaking,
        when she heard suddenly some one outside say, "Miss Lin is come;" and
        shortly Lin Tai-yü walked in in a jaunty manner.}
}
{
}

\Verse{38}
{
    \zA{}
}
{
    \eA{"Oh, I come at a wrong moment!" she exclaimed forthwith, smirking
        significantly when she caught sight of Pao-yü. Pao-yü and the rest lost no
        time in rising and offering her a seat, whereupon Pao Ch'ai added with a
        smile, "How can you say such things?" "Had I known sooner," continued
        Tai-yü, "that he was here, I would have kept away." "I can't fathom this
        meaning of yours," protested Pao Ch'ai. "If one comes," Tai-yü urged
        smiling, "then all come, and when one doesn't come, then no one comes. Now
        were he to come to-day, and I to come to-morrow, wouldn't there be, by a
        division of this kind, always some one with you every day? and in this way,
        you wouldn't feel too lonely, nor too crowded. How is it, cousin, that you
        didn't understand what I meant to imply?" "Is it snowing?"}
}
{
}

\Verse{39}
{
    \zA{}
}
{
    \eA{inquired Pao-yü, upon noticing that she wore a cloak made of crimson camlet,
        buttoning in front. "It has been snowing for some time," ventured the
        matrons, who were standing below. "Fetch my wrapper!" Pao-yü remarked, and
        Tai-yü readily laughed. "Am I not right? I come, and, of course, he must go
        at once." "Did I ever mention that I was going?" questioned Pao-yü; "I only
        wish it brought to have it ready when I want it." "It's a snowy day,"
        consequently remarked Pao-yü's nurse, dame Li, "and we must also look to the
        time, but you had better remain here and amuse yourself with your cousin.
        Your aunt has, in there, got ready tea and fruits. I'll tell the
        waiting-maid to go and fetch your wrapper and the boys to return home."}
}
{
}

\Verse{40}
{
    \zA{}
}
{
    \eA{Pao-yü assented, and nurse Li left the room and told the boys that they were
        at liberty to go. By this time Mrs. Hsüeh had prepared tea and several kinds
        of nice things and kept them all to partake of those delicacies. Pao-yü,
        having spoken highly of some goose feet and ducks' tongues he had tasted
        some days before, at his eldest sister-in-law's, Mrs. Yu's, "aunt" Hsüeh
        promptly produced several dishes of the same kind, made by herself, and gave
        them to Pao-yü to try. "With a little wine," added Pao-yü with a smile,
        "they would be first rate." Mrs. Hsüeh thereupon bade the servants fetch
        some wine of the best quality;}
}
{
}

\Verse{41}
{
    \zA{}
}
{
    \eA{but dame Li came forward and remonstrated. "My lady," she said, "never mind
        the wine." Pao-yü smilingly pleaded: "My nurse, I'll take just one cup and
        no more." "It's no use," nurse Li replied, "were your grandmother and mother
        present, I wouldn't care if you drank a whole jar. I remember the day when I
        turned my eyes away but for a moment, and some ignorant fool or other,
        merely with the view of pandering for your favour, gave you only a drop of
        wine to drink, and how this brought reproaches upon me for a couple of days.
        You don't know, my lady, you have no idea of his disposition! it's really
        dreadful; and when he has had a little wine he shows far more temper.}
}
{
}

\Verse{42}
{
    \zA{}
}
{
    \eA{On days when her venerable ladyship is in high spirits, she allows him to
        have his own way about drinking, but he's not allowed to have wine on any
        and every day; and why should I have to suffer inside and all for nothing at
        all?" "You antiquated thing!" replied Mrs. Hsüeh laughing, "set your mind at
        ease, and go and drink your own wine! I won't let him have too much, and
        should even the old lady say anything, let the fault be mine." Saying this,
        she asked a waiting-maid to take nurse Li along with her and give her also a
        glass of wine so as to keep out the cold air. When nurse Li heard these
        words, she had no alternative but to go for a time with all the others and
        have some wine to drink. "The wine need not be warmed: I prefer it cold!"}
}
{
}

\Verse{43}
{
    \zA{}
}
{
    \eA{Pao-yü went on to suggest meanwhile. "That won't do," remonstrated Mrs.
        Hsüeh; "cold wine will make your hand tremble when you write." "You have,"
        interposed Pao Ch'ai smiling, "the good fortune, cousin Pao-yü, of having
        daily opportunities of acquiring a knowledge of every kind of subject, and
        yet don't you know that the properties of wine are mostly heating? If you
        drink wine warm, its effects soon dispel, but if you drink it cold, it at
        once congeals in you; and as upon your intestines devolves the warming of
        it, how can you not derive any harm? and won't you yet from this time change
        this habit of yours? leave off at once drinking that cold wine."}
}
{
}

\Verse{44}
{
    \zA{}
}
{
    \eA{Pao-yü finding that the words he had heard contained a good deal of sense,
        speedily put down the cold wine, and having asked them to warm it, he at
        length drank it. Tai-yü was bent upon cracking melon seeds, saying nothing
        but simply pursing up her lips and smiling, when, strange coincidence, Hsüeh
        Yen, Tai-yü's waiting-maid, walked in and handed her mistress a small
        hand-stove. "Who told you to bring it?" ascertained Tai-yü grinningly. "I'm
        sorry to have given whoever it is the trouble; I'm obliged to her. But did
        she ever imagine that I would freeze to death?" "Tzu Chuan was afraid,"
        replied Hsüeh Yen, "that you would, miss, feel cold, and she asked me to
        bring it over." Tai-yü took it over and held it in her lap.}
}
{
}

\Verse{45}
{
    \zA{}
}
{
    \eA{"How is it," she smiled, "that you listen to what she tells you, but that
        you treat what I say, day after day, as so much wind blowing past your ears!
        How is it that you at once do what she bids you, with even greater alacrity
        than you would an imperial edict?" When Pao-yü heard this, he felt sure in
        his mind that Tai-yü was availing herself of this opportunity to make fun of
        him, but he made no remark, merely laughing to himself and paying no further
        notice. Pao Ch'ai, again, knew full well that this habit was a weak point
        with Tai-yü, so she too did not go out of her way to heed what she said.}
}
{
}

\Verse{46}
{
    \zA{}
}
{
    \eA{"You've always been delicate and unable to stand the cold," interposed
        "aunt" Hsüeh, "and is it not a kind attention on their part to have thought
        of you?" "You don't know, aunt, how it really stands," responded Tai-yü
        smilingly; "fortunately enough, it was sent to me here at your quarters; for
        had it been in any one else's house, wouldn't it have been a slight upon
        them? Is it forsooth nice to think that people haven't so much as a
        hand-stove, and that one has fussily to be sent over from home? People won't
        say that the waiting-maids are too officious, but will imagine that I'm in
        the habit of behaving in this offensive fashion." "You're far too
        punctilious," remarked Mrs. Hsüeh, "as to entertain such notions! No such
        ideas as these crossed my mind just now."}
}
{
}

\Verse{47}
{
    \zA{}
}
{
    \eA{While they were conversing, Pao-yü had taken so much as three cups of wine,
        and nurse Li came forward again to prevent him from having any more. Pao-yü
        was just then in a state of exultation and excitement, (a state) enhanced by
        the conversation and laughter of his cousins, so that was he ready to agree
        to having no more! But he was constrained in a humble spirit to entreat for
        permission. "My dear nurse," he implored, "I'll just take two more cups and
        then have no more." "You'd better be careful," added nurse Li, "your father
        is at home to-day, and see that you're ready to be examined in your
        lessons." When Pao-yü heard this mention, his spirits at once sank within
        him, and gently putting the wine aside, he dropped his head upon his breast.
        Tai-yü promptly remonstrated.}
}
{
}

\Verse{48}
{
    \zA{}
}
{
    \eA{"You've thrown cold water," she said, "over the spirits of the whole
        company; why, if uncle should ask to see you, well, say that aunt Hsüeh
        detained you. This old nurse of yours has been drinking, and again makes us
        the means of clearing her muddled head!" While saying this, she gave Pao-yü
        a big nudge with the intent of stirring up his spirits, adding, as she
        addressed him in a low tone of voice: "Don't let us heed that old thing, but
        mind our own enjoyment." Dame Li also knew very well Tai-yü's disposition,
        and therefore remarked: "Now, Miss Lin, don't you urge him on;}
}
{
}

\Verse{49}
{
    \zA{}
}
{
    \eA{you should after all, give him good advice, as he may, I think, listen to a
        good deal of what you say to him." "Why should I urge him on?" rejoined Lin
        Tai-yü, with a sarcastic smile, "nor will I trouble myself to give him
        advice. You, old lady, are far too scrupulous! Old lady Chia has also time
        after time given him wine, and if he now takes a cup or two more here, at
        his aunt's, lady Hsüeh's house, there's no harm that I can see. Is it
        perhaps, who knows, that aunt is a stranger in this establishment, and that
        we have in fact no right to come over here to see her?" Nurse Li was both
        vexed and amused by the words she had just heard.}
}
{
}

\Verse{50}
{
    \zA{}
}
{
    \eA{"Really," she observed, "every remark this girl Lin utters is sharper than a
        razor! I didn't say anything much!" Pao Ch'ai too could not suppress a
        smile, and as she pinched Tai-yü's cheek, she exclaimed, "Oh the tongue of
        this frowning girl! one can neither resent what it says, nor yet listen to
        it with any gratification!" "Don't be afraid!" Mrs. Hsüeh went on to say,
        "don't be afraid; my son, you've come to see me, and although I've nothing
        good to give you, you mustn't, through fright, let the trifle you've taken
        lie heavy on your stomach, and thus make me uneasy; but just drink at your
        pleasure, and as much as you like, and let the blame fall on my shoulders.
        What's more, you can stay to dinner with me, and then go home; or if you do
        get tipsy, you can sleep with me, that's all."}
}
{
}

\Verse{51}
{
    \zA{}
}
{
    \eA{She thereupon told the servants to heat some more wine. "I'll come," she
        continued, "and keep you company while you have two or three cups, after
        which we'll have something to eat!" It was only after these assurances that
        Pao-yü's spirits began at length, once more to revive, and dame Li then
        directed the waiting-maids what to do. "You remain here," she enjoined, "and
        mind, be diligent while I go home and change; when I'll come back again.
        Don't allow him," she also whispered to "aunt" Hsüeh, "to have all his own
        way and drink too much." Having said this, she betook herself back to her
        quarters; and during this while, though there were two or three nurses in
        attendance, they did not concern themselves with what was going on.}
}
{
}

\Verse{52}
{
    \zA{}
}
{
    \eA{As soon as they saw that nurse Li had left, they likewise all quietly
        slipped out, at the first opportunity they found, while there remained but
        two waiting-maids, who were only too glad to curry favour with Pao-yü. But
        fortunately "aunt" Hsüeh, by much coaxing and persuading, only let him have
        a few cups, and the wine being then promptly cleared away, pickled bamboo
        shoots and chicken-skin soup were prepared, of which Pao-yü drank with
        relish several bowls full, eating besides more than half a bowl of finest
        rice congee. By this time, Hsüeh Pao Ch'ai and Lin Tai-yü had also finished
        their repast;}
}
{
}

\Verse{53}
{
    \zA{}
}
{
    \eA{and when Pao-yü had drunk a few cups of strong tea, Mrs. Hsüeh felt more
        easy in her mind. Hsüeh Yen and the others, three or four of them in all,
        had also had their meal, and came in to wait upon them. "Are you now going
        or not?" inquired Tai-yü of Pao-yü. Pao-yü looked askance with his drowsy
        eyes. "If you want to go," he observed, "I'll go with you." Tai-yü hearing
        this, speedily rose. "We've been here nearly the whole day," she said, "and
        ought to be going back." As she spoke the two of them bade good-bye, and the
        waiting-maids at once presented a hood to each of them. Pao-yü readily
        lowered his head slightly and told a waiting-maid to put it on. The girl
        promptly took the hood, made of deep red cloth, and shaking it out of its
        folds, she put it on Pao-yü's head.}
}
{
}

\Verse{54}
{
    \zA{}
}
{
    \eA{"That will do," hastily exclaimed Pao-yü. "You stupid thing! gently a bit;
        is it likely you've never seen any one put one on before? let me do it
        myself." "Come over here, and I'll put it on for you," suggested Tai-yü, as
        she stood on the edge of the couch. Pao-yü eagerly approached her, and
        Tai-yü carefully kept the cap, to which his hair was bound, fast down, and
        taking the hood she rested its edge on the circlet round his forehead. She
        then raised the ball of crimson velvet, which was as large as a walnut, and
        put it in such a way that, as it waved tremulously, it should appear outside
        the hood.}
}
{
}

\Verse{55}
{
    \zA{}
}
{
    \eA{These arrangements completed she cast a look for a while at what she had
        done. "That's right now," she added, "throw your wrapper over you!" When
        Pao-yü caught these words, he eventually took the wrapper and threw it over
        his shoulders. "None of your nurses," hurriedly interposed aunt Hsüeh, "are
        yet come, so you had better wait a while." "Why should we wait for them?"
        observed Pao-yü. "We have the waiting-maids to escort us, and surely they
        should be enough." Mrs. Hsüeh finding it difficult to set her mind at ease
        deputed two married women to accompany the two cousins; and after they had
        both expressed (to these women) their regret at having troubled them, they
        came straightway to dowager lady Chia's suite of apartments. Her venerable
        ladyship had not, as yet, had her evening repast.}
}
{
}

\Verse{56}
{
    \zA{}
}
{
    \eA{Hearing that they had been at Mrs. Hsüeh's, she was extremely pleased; but
        noticing that Pao-yü had had some wine, she gave orders that he should be
        taken to his room, and put to bed, and not be allowed to come out again. "Do
        take good care of him," she therefore enjoined the servants, and when
        suddenly she bethought herself of Pao-yü's attendants, "How is it," she at
        once inquired of them all, "that I don't see nurse Li here?" They did not
        venture to tell her the truth, that she had gone home, but simply explained
        that she had come in a few moments back, and that they thought she must have
        again gone out on some business or other.}
}
{
}

\Verse{57}
{
    \zA{}
}
{
    \eA{"She's better off than your venerable ladyship," remarked Pao-yü, turning
        round and swaying from side to side. "Why then ask after her? Were I rid of
        her, I believe I might live a little longer." While uttering these words, he
        reached the door of his bedroom, where he saw pen and ink laid out on the
        writing table. "That's nice," exclaimed Ch'ing Wen, as she came to meet him
        with a smile on her face, "you tell me to prepare the ink for you, but
        though when you get up, you were full of the idea of writing, you only wrote
        three characters, when you discarded the pencil, and ran away, fooling me,
        by making me wait the whole day! Come now at once and exhaust all this ink
        before you're let off."}
}
{
}

\Verse{58}
{
    \zA{}
}
{
    \eA{Pao-yü then remembered what had taken place in the morning. "Where are the
        three characters I wrote?" he consequently inquired, smiling. "Why this man
        is tipsy," remarked Ch'ing Wen sneeringly. "As you were going to the other
        mansion, you told me to stick them over the door. I was afraid lest any one
        else should spoil them, as they were being pasted, so I climbed up a high
        ladder and was ever so long in putting them up myself; my hands are even now
        numb with cold." "Oh I forgot all about it," replied Pao-yü grinning, "if
        your hands are cold, come and I'll rub them warm for you."}
}
{
}

\Verse{59}
{
    \zA{}
}
{
    \eA{Promptly stretching out his hand, he took those of Ch'ing Wen in his, and
        the two of them looked at the three characters, which he recently had
        written, and which were pasted above the door. In a short while, Tai-yü
        came. "My dear cousin," Pao-yü said to her smilingly, "tell me without any
        prevarication which of the three characters is the best written?" Tai-yü
        raised her head and perceived the three characters: Red, Rue, Hall. "They're
        all well done," she rejoined, with a smirk, "How is it you've written them
        so well? By and bye you must also write a tablet for me." "Are you again
        making fun of me?" asked Pao-yü smiling; "what about sister Hsi Jen?" he
        went on to inquire.}
}
{
}

\Verse{60}
{
    \zA{}
}
{
    \eA{Ch'ing Wen pouted her lips, pointing towards the stove-couch in the inner
        room, and, on looking in, Pao-yü espied Hsi Jen fast asleep in her daily
        costume. "Well," Pao-yü observed laughing, "there's no harm in it, but its
        rather early to sleep. When I was having my early meal, on the other side,"
        he proceeded, speaking to Ch'ing Wen, "there was a small dish of dumplings,
        with bean-curd outside; and as I thought you would like to have some, I
        asked Mrs. Yu for them, telling her that I would keep them, and eat them in
        the evening; I told some one to bring them over, but have you perchance seen
        them?" "Be quick and drop that subject," suggested Ch'ing Wen; "as soon as
        they were brought over, I at once knew they were intended for me;}
}
{
}

\Verse{61}
{
    \zA{}
}
{
    \eA{as I had just finished my meal, I put them by in there, but when nurse Li
        came she saw them. 'Pao-yü,' she said, 'is not likely to eat them, so I'll
        take them and give them to my grandson.' And forthwith she bade some one
        take them over to her home." While she was speaking, Hsi Hsüeh brought in
        tea, and Pao-yü pressed his cousin Lin to have a cup. "Miss Lin has gone
        long ago," observed all of them, as they burst out laughing, "and do you
        offer her tea?" Pao-yü drank about half a cup, when he also suddenly
        bethought himself of some tea, which had been brewed in the morning. "This
        morning," he therefore inquired of Hsi Hsüeh, "when you made a cup of
        maple-dew tea, I told you that that kind of tea requires brewing three or
        four times before its colour appears;}
}
{
}

\Verse{62}
{
    \zA{}
}
{
    \eA{and how is that you now again bring me this tea?" "I did really put it by,"
        answered Hsi Hsüeh, "but nurse Li came and drank it, and then went off."
        Pao-yü upon hearing this, dashed the cup he held in his hand on the ground,
        and as it broke into small fragments, with a crash, it spattered Hsi Hsüeh's
        petticoat all over. "Of whose family is she the mistress?" inquired Pao-yü
        of Hsi Hsüeh, as he jumped up, "that you all pay such deference to her. I
        just simply had a little of her milk, when I was a brat, and that's all; and
        now she has got into the way of thinking herself more high and mighty than
        even the heads of the family! She should be packed off, and then we shall
        all have peace and quiet."}
}
{
}

\Verse{63}
{
    \zA{}
}
{
    \eA{Saying this, he was bent upon going, there and then, to tell dowager lady
        Chia to have his nurse driven away. Hsi Jen was really not asleep, but
        simply feigning, with the idea, when Pao-yü came, to startle him in play. At
        first, when she heard him speak of writing, and inquire after the dumplings,
        she did not think it necessary to get up, but when he flung the tea-cup on
        the floor, and got into a temper, she promptly jumped up and tried to
        appease him, and to prevent him by coaxing from carrying out his threat. A
        waiting-maid sent by dowager lady Chia came in, meanwhile, to ask what was
        the matter.}
}
{
}

\Verse{64}
{
    \zA{}
}
{
    \eA{"I had just gone to pour tea," replied Hsi Jen, without the least
        hesitation, "and I slipped on the snow and fell, while the cup dropped from
        my hand and broke. Your decision to send her away is good," she went on to
        advise Pao-yü, "and we are all willing to go also; and why not avail
        yourself of this opportunity to dismiss us in a body? It will be for our
        good, and you too on the other hand, needn't perplex yourself about not
        getting better people to come and wait on you!" When Pao-yü heard this
        taunt, he had at length not a word to say, and supported by Hsi Jen and the
        other attendants on to the couch, they divested him of his clothes.}
}
{
}

\Verse{65}
{
    \zA{}
}
{
    \eA{But they failed to understand the drift of what Pao-yü kept on still
        muttering, and all they could make out was an endless string of words; but
        his eyes grew heavier and drowsier, and they forthwith waited upon him until
        he went to sleep; when Hsi Jen unclasped the jade of spiritual perception,
        and rolling it up in a handkerchief, she lay it under the mattress, with the
        idea that when he put it on the next day it should not chill his neck.
        Pao-yü fell sound asleep the moment he lay his head on the pillow.}
}
{
}

\Verse{66}
{
    \zA{}
}
{
    \eA{By this time nurse Li and the others had come in, but when they heard that
        Pao-yü was tipsy, they too did not venture to approach, but gently made
        inquiries as to whether he was asleep or not. On hearing that he was, they
        took their departure with their minds more at ease. The next morning the
        moment Pao-yü awoke, some one came in to tell him that young Mr. Jung,
        living in the mansion on the other side, had brought Ch'in Chung to pay him
        a visit. Pao-yü speedily went out to greet them and to take them over to pay
        their respects to dowager lady Chia.}
}
{
}

\Verse{67}
{
    \zA{}
}
{
    \eA{Her venerable ladyship upon perceiving that Ch'in Chung, with his handsome
        countenance, and his refined manners, would be a fit companion for Pao-yü in
        his studies, felt extremely delighted at heart; and having readily detained
        him to tea, and kept him to dinner, she went further and directed a servant
        to escort him to see madame Wang and the rest of the family. With the fond
        regard of the whole household for Mrs. Ch'in, they were, when they saw what
        a kind of person Ch'in Chung was, so enchanted with him, that at the time of
        his departure, they all had presents to give him; even dowager lady Chia
        herself presented him with a purse and a golden image of the God of
        Learning, with a view that it should incite him to study and harmony.}
}
{
}

\Verse{68}
{
    \zA{}
}
{
    \eA{"Your house," she further advised him, "is far off, and when it's cold or
        hot, it would be inconvenient for you to come all that way, so you had
        better come and live over here with me. You'll then be always with your
        cousin Pao-yü, and you won't be together, in your studies, with those
        fellow-pupils of yours who have no idea what progress means." Ch'in Chung
        made a suitable answer to each one of her remarks, and on his return home he
        told everything to his father. His father, Ch'in Pang-yeh, held at present
        the post of Secretary in the Peking Field Force, and was well-nigh seventy.
        His wife had died at an early period, and as she left no issue, he adopted a
        son and a daughter from a foundling asylum.}
}
{
}

\Verse{69}
{
    \zA{}
}
{
    \eA{But who would have thought it, the boy also died, and there only remained
        the girl, known as Kó Ch'ing in her infancy, who when she grew up, was
        beautiful in face and graceful in manners, and who by reason of some
        relationship with the Chia family, was consequently united by the ties of
        marriage (to one of the household). Ch'in Pang-yeh was in his fiftieth year
        when he at length got this son. As his tutor had the previous year left to
        go south, he remained at home keeping up his former lessons; and (his
        father) had been just thinking of talking over the matter with his relatives
        of the Chia family, and sending his son to the private school, when, as luck
        would have it, this opportunity of meeting Pao-yü presented itself.}
}
{
}

\Verse{70}
{
    \zA{}
}
{
    \eA{Knowing besides that the family school was under the direction of the
        venerable scholar Chia Tai-ju, and hoping that by joining his class, (his
        son) might advance in knowledge and by these means reap reputation, he was
        therefore intensely gratified. The only drawbacks were that his official
        emoluments were scanty, and that both the eyes of everyone in the other
        establishment were set upon riches and honours, so that he could not
        contribute anything short of the amount (given by others); but his son's
        welfare throughout life was a serious consideration, and he, needless to
        say, had to scrape together from the East and to collect from the West;}
}
{
}

\Verse{71}
{
    \zA{}
}
{
    \eA{and making a parcel, with all deference, of twenty-four taels for an
        introduction present, he came along with Ch'in Chung to Tai-ju's house to
        pay their respects. But he had to wait subsequently until Pao-yü could fix
        on an auspicious date on which they could together enter the school. As for
        what happened after they came to school, the next chapter will divulge.}
}
{
}
